% from ega0-20-fr.tex
\oldpage[0\textsubscript{IV-1}]{128}
\begin{env}[20.5.1]
\label{0.20.5.1-fr}
Dans tout ce numéro et le suivant, sauf mention expresse du contraire, les
anneaux et modules considérés sont supposés munis de la topologie
\emph{discrète}.
\end{env}

\begin{env}[20.5.2]
\label{0.20.5.2-fr}
Soient $A$ un anneau, $B$, $C$ deux $A$-algèbres, $u:B\to C$ un
$A$-homomorphisme ; on a un diagramme commutatif
\[
\xymatrix{
  B\otimes_A B \ar[r]^{u\otimes u} \ar[d]_{p_{B/A}} &
  C\otimes_A C \ar[d]^{p_{C/A}} \\
  B \ar[r]^u & C
}
\tag{20.5.2.1}
\label{0.20.5.2.1-fr}
\]
d'où par passage aux quotients, un $A$-homomorphisme d'algèbres
\begin{equation}
\tag{20.5.2.2}
\label{0.20.5.2.2-fr}
  u':\mathrm{P}^1_{B/A}\to\mathrm{P}^1_{C/A}
\end{equation}
tel que le diagramme
\[
\xymatrix{
  \mathrm{P}^1_{B/A} \ar[r]^{u'} & \mathrm{P}^1_{C/A} \\
  B \ar[u]^{p_1} \ar[r]^u & C \ar[u]_{p_1}
}
\]
soit commutatif ; comme $u\otimes u$ applique $\mathfrak{J}_{B/A}$ dans
$\mathfrak{J}_{C/A}$, on obtient, en restreignant $u'$ à
$\Omega^1_{B/A}$, une application
\begin{equation}
\tag{20.5.2.3}
\label{0.20.5.2.3-fr}
  u'':\Omega^1_{B/A}\to\Omega^1_{C/A}
\end{equation}
telle que le couple $(u'',u)$ soit un \emph{di-homomorphisme} pour la
structure de $B$-module de $\Omega^1_{B/A}$ et la structure de $C$-module de
$\Omega^1_{C/A}$ ; ce dernier fait permet d'en déduire canoniquement un
\emph{homomorphisme de $C$-modules}
\begin{equation}
\tag{20.5.2.4}
\label{0.20.5.2.4-fr}
  u_{C/B/A}:\Omega^1_{B/A}\otimes_B C\to\Omega^1_{C/A}.
\end{equation}

